% !PRINTED 85
\weberpage{85}
%<<CONT
determine the factor $Q$. One obtains first
\[
1 = Q^2\prod^{\nu}\frac{(1-q^{2\nu-1})^2(1+q^{2\nu-1})^2(1+q^{2\nu})^2}{(1-q^{2\nu})^2},
\]
or, by disposing of the still undetermined sign of $Q$ and thereby of the \gesperrt{signs} of the
$\vartheta$-functions:
\[
Q = \prod_{1,\infty}^{\nu}\frac{(1-q^{2\nu})}{(1+q^{2\nu})(1+q^{2\nu-1})(1-q^{2\nu-1})}.
\]
In the denominator one may put $\prod(1+q^{\nu})$ in place of $\prod(1+q^{2\nu})(1+q^{2\nu-1})$,
and the numerator $\prod(1-q^{2\nu})$ can be decomposed into
\[
\prod(1-q^{\nu})(1+q^{\nu}) = \prod(1-q^{2\nu})(1-q^{2\nu-1})(1+q^{\nu}).
\]
In this way there results finally
\weq{6}{Q = \prod^{\nu}(1 - q^{2\nu}).}
The expressions (1), (2), (3), (4) can be represented in real form if one carries out the
multiplication of the conjugate imaginary factors. One thus obtains:
\weq{7}{\begin{gathered}
\vartheta_{00}(v) = \prod_{1,\infty}^{\nu}(1 - q^{2\nu})(1 + 2q^{2\nu-1}\cos 2\pi v + q^{4\nu-2}),\\
\vartheta_{01}(v) = \prod_{1,\infty}^{\nu}(1 - q^{2\nu})(1 - 2q^{2\nu-1}\cos 2\pi v + q^{4\nu-1}),\\
% [sic] German source prints exponent 4\nu-1; mathematically expected 4\nu-2.
\vartheta_{10}(v) = 2q^{\frac{1}{4}}\cos\pi v\prod_{1,\infty}^{\nu}(1 - q^{2\nu})(1 + 2q^{2\nu}\cos 2\pi v + q^{4\nu}),\\
\vartheta_{11}(v) = 2q^{\frac{1}{4}}\sin\pi v\prod_{1,\infty}^{\nu}(1 - q^{2\nu})(1 - 2q^{2\nu}\cos 2\pi v + q^{4\nu}).
\end{gathered}}
If we introduce the expression (6) into the last equation (5), there results
\[
\vartheta_{11}' = 2\pi q^{\frac{1}{4}}\prod(1 - q^{2\nu})^3,
\]
and if we therefore put
\weq{8}{\eta(\omega) = q^{\frac{1}{12}}\prod(1 - q^{2\nu})}
then we have
\weq{9}{\vartheta_{11}' = 2\pi\eta(\omega)^3.}
By eliminating $Q$ with the help of the relation
\[
Q = q^{-\frac{1}{12}}\eta(\omega)
\]
from (5), we set
\weq{10}{\begin{gathered}
\vartheta_{00} = \eta(\omega)f(\omega)^2,\\
\vartheta_{01} = \eta(\omega)f_1(\omega)^2,\\
\vartheta_{10} = \eta(\omega)f_2(\omega)^2,
\end{gathered}}
%CONT>>
